\documentclass[11pt]{article}

\usepackage[margin=1in]{geometry}
\usepackage{amsmath,amssymb,bm}
\usepackage{enumitem}
\usepackage{hyperref}

\newcommand{\Ga}{\mathrm{Ga}}
\newcommand{\IGa}{\mathrm{IGa}}
\newcommand{\Cat}{\mathrm{Cat}}
\newcommand{\Dir}{\mathrm{Dirichlet}}
\newcommand{\PIG}{\mathrm{P\mbox{-}IG}}
\newcommand{\PTN}{\mathrm{PTN}}
\newcommand{\diff}{\,\mathrm{d}}
\newcommand{\E}{\mathbb{E}}

\title{P-IG MCMC Sampler for the Original GaGa Model}
\author{}
\date{}

\begin{document}
\maketitle

\section{Target GaGa Model}

This document writes the P-IG data-augmentation MCMC sampler for the
GaGa model as implemented in Rossell's \texttt{gaga} package.  The main
version below is the default \texttt{equalcv=TRUE, nclust=1} GaGa model,
which is the version used in the Armstrong and MAQC benchmarks.  The notation
matches the package parameters
\[
  a_0=\texttt{a0}, \qquad \nu=\texttt{nu}, \qquad
  \beta=\texttt{balpha}, \qquad \mu=\texttt{nualpha}.
\]
All gamma distributions use the shape-rate convention.

Let \(x_{ij}\) be the expression of gene \(i=1,\ldots,G\) in array
\(j=1,\ldots,n\), and let \(g_j\in\{1,\ldots,K\}\) be its experimental group.
Let \(h=0,\ldots,H-1\) index expression patterns.  Pattern \(h\) maps groups
to clusters through
\[
  c_h(g)\in\{1,\ldots,C_h\}.
\]
For example, the null pattern has \(C_h=1\), while the full alternative has
one cluster per group.

The default constant-CV GaGa model is
\begin{align}
  x_{ij}\mid \alpha_i,\lambda_{i,c_h(g_j)},\delta_i=h
    &\sim \Ga\!\left(\alpha_i,\,
      \frac{\alpha_i}{\lambda_{i,c_h(g_j)}}\right),
      \label{eq:obs}\\
  \lambda_{ic}\mid a_0,\nu
    &\sim \IGa\!\left(a_0,\frac{a_0}{\nu}\right),
      \qquad c=1,\ldots,C_{\delta_i}, \label{eq:lambda-prior}\\
  \alpha_i\mid \beta,\mu
    &\sim \Ga\!\left(\beta,\frac{\beta}{\mu}\right),
      \label{eq:alpha-prior}\\
  \delta_i\mid \bm\pi
    &\sim \Cat(\pi_0,\ldots,\pi_{H-1}). \label{eq:pattern-prior}
\end{align}
Equivalently, with \(q_{ic}=1/\lambda_{ic}\),
\[
  q_{ic}\mid a_0,\nu \sim \Ga\!\left(a_0,\frac{a_0}{\nu}\right).
\]

The hyperpriors used by the fully Bayesian GaGa model can be written as
\begin{align}
  a_0 &\sim \Ga(a_{a_0},b_{a_0}), &
  \nu &\sim \IGa(a_\nu,b_\nu), \nonumber\\
  \beta &\sim \Ga(a_\beta,b_\beta), &
  \mu &\sim \IGa(a_\mu,b_\mu), \nonumber\\
  \bm\pi &\sim \Dir(p_0,\ldots,p_{H-1}). \label{eq:hyperpriors}
\end{align}

\section{Sufficient Statistics Under a Pattern}

For gene \(i\), pattern \(h\), and cluster \(c\), define
\begin{align}
  I_{hc} &= \{j: c_h(g_j)=c\}, &
  n_{hc} &= |I_{hc}|, \nonumber\\
  S_{ihc} &= \sum_{j\in I_{hc}} x_{ij}, &
  L_{ihc} &= \sum_{j\in I_{hc}} \log x_{ij}, \nonumber\\
  r_0 &= \frac{a_0}{\nu}.
\end{align}

After integrating \(q_{ic}=1/\lambda_{ic}\), the cluster marginal likelihood
conditional on \(\alpha_i=\alpha\) is
\begin{align}
  m_{ihc}(\alpha)
  &=
  \int p(\bm x_{ihc}\mid \alpha,q)\,p(q\mid a_0,\nu)\,\diff q \nonumber\\
  &=
  \frac{r_0^{a_0}}{\Gamma(a_0)}
  \frac{\Gamma(a_0+n_{hc}\alpha)}{\Gamma(\alpha)^{n_{hc}}}
  \frac{
    \alpha^{n_{hc}\alpha}
    \exp\{(\alpha-1)L_{ihc}\}
  }{
    (r_0+\alpha S_{ihc})^{a_0+n_{hc}\alpha}
  }.
  \label{eq:cluster-marginal}
\end{align}
The pattern likelihood is
\[
  m_{ih}(\alpha)=\prod_{c=1}^{C_h} m_{ihc}(\alpha).
\]

\section{Pattern Update}

Given \(\alpha_i\) and the hyperparameters, update the expression pattern by
\begin{align}
  \Pr(\delta_i=h\mid \alpha_i,\bm x_i,\Theta)
  =
  \frac{\pi_h m_{ih}(\alpha_i)}
       {\sum_{\ell=0}^{H-1}\pi_\ell m_{i\ell}(\alpha_i)}.
  \label{eq:pattern-update}
\end{align}
This is the same pattern posterior structure used by \texttt{gaga}; the
difference is that \(\alpha_i\) is sampled with P-IG augmentation rather than
using Rossell's gamma-shape approximation.

\section{Collapsed Conditional for \texorpdfstring{\(\alpha_i\)}{alpha}}

Given \(\delta_i=h\), the collapsed conditional density of \(\alpha_i\) is
proportional to
\begin{align}
  p(\alpha_i=\alpha \mid \delta_i=h,\bm x_i,\Theta)
  &\propto
  \alpha^{\beta-1}\exp\!\left(-\frac{\beta}{\mu}\alpha\right)
  \prod_{c=1}^{C_h} m_{ihc}(\alpha).
  \label{eq:alpha-target}
\end{align}
Dropping constants that do not depend on \(\alpha\),
\begin{align}
  \log p(\alpha\mid\cdot)
  &=
  (\beta-1)\log\alpha-\frac{\beta}{\mu}\alpha \nonumber\\
  &\quad+
  \sum_{c=1}^{C_h}
  \left[
    n_{hc}\alpha\log\alpha
    -n_{hc}\log\Gamma(\alpha)
    +(\alpha-1)L_{ihc}
    +\log\Gamma(a_0+n_{hc}\alpha)
    -(a_0+n_{hc}\alpha)\log(r_0+\alpha S_{ihc})
  \right].
  \label{eq:alpha-log-target}
\end{align}

\section{P-IG Augmentation for the Gamma-Shape Conditional}

Let
\[
  N_h=\sum_{c=1}^{C_h} n_{hc}.
\]
The reciprocal gamma terms are handled by the P-IG identity
\[
  \frac{1}{\Gamma(\alpha)^{N_h}}
  \propto
  \alpha^{N_h}\exp(N_h\gamma\alpha)\,
  \E\!\left[\exp\!\left(-\alpha^2\sum_{s=1}^{N_h}\omega_s\right)\right],
  \qquad
  \omega_s\mid \alpha\sim \PIG(\alpha),
\]
where \(\gamma=-\psi(1)\) is the Euler--Mascheroni constant.

The numerator gamma terms are augmented by
\[
  \Gamma(a_0+n_{hc}\alpha)
  =
  \int_0^\infty
  \tau_c^{a_0+n_{hc}\alpha-1}\exp(-\tau_c)\,\diff\tau_c,
  \qquad
  \tau_c\mid\alpha\sim \Ga(a_0+n_{hc}\alpha,1).
\]

Given the auxiliary variables, define
\[
  \Omega=\sum_{s=1}^{N_h}\omega_s
\]
and
\begin{align}
  B_{ih}
  =
  N_h\gamma
  +\sum_{c=1}^{C_h} L_{ihc}
  -\frac{\beta}{\mu}
  +\sum_{c=1}^{C_h} n_{hc}\log\tau_c.
  \label{eq:Bih}
\end{align}
The augmented conditional for \(\alpha\) has the form
\begin{align}
  p(\alpha\mid\bm\tau,\bm\omega,\delta_i=h,\bm x_i,\Theta)
  \propto
  \alpha^{N_h+\beta-1}
  \exp\{-\Omega\alpha^2+B_{ih}\alpha\}
  \exp\{R_{ih}(\alpha)\},
  \label{eq:aug-alpha}
\end{align}
where the non-PTN remainder is
\begin{align}
  R_{ih}(\alpha)
  =
  N_h\alpha\log\alpha
  -
  \sum_{c=1}^{C_h}
  (a_0+n_{hc}\alpha)\log(r_0+\alpha S_{ihc}).
  \label{eq:remainder}
\end{align}

\section{Exact Log-Scale Slice Step for \texorpdfstring{\(\alpha_i\)}{alpha}}

The benchmark implementation samples the augmented conditional
\eqref{eq:aug-alpha} directly on the log scale.  Let \(y=\log\alpha\).  After
including the Jacobian, the log density for \(y\) is
\begin{align}
  \ell_{ih}(y)
  =
  (N_h+\beta)y
  -\Omega\exp(2y)
  +B_{ih}\exp(y)
  +R_{ih}\{\exp(y)\}.
  \label{eq:log-alpha-slice}
\end{align}
The code uses a standard stepping-out and shrinkage slice sampler for
\(\ell_{ih}(y)\), then sets \(\alpha_i=\exp(y)\).  This is an exact
one-dimensional MCMC update for the same P-IG-augmented conditional, and it
avoids the near-zero acceptance of the plain PTN independence proposal in the
real microarray benchmarks.

\paragraph{Optional PTN independence proposal.}
One exact diagnostic proposal uses only the PTN part of \eqref{eq:aug-alpha}:
\[
  \alpha^\star
  \sim
  \PTN\!\left(
    p=N_h+\beta,\,
    a=\Omega,\,
    b=B_{ih}
  \right),
\]
where
\[
  \PTN(p,a,b):\quad f(x)\propto x^{p-1}\exp(-a x^2+bx),\qquad x>0.
\]
Since this proposal is independent of the current \(\alpha\) conditional on
\((\bm\tau,\bm\omega)\), the PTN terms cancel in the Metropolis ratio, and
\[
  A(\alpha\to\alpha^\star)
  =
  \min\left\{
    1,\,
    \exp\!\left[
      R_{ih}(\alpha^\star)-R_{ih}(\alpha)
    \right]
  \right\}.
  \label{eq:exact-alpha-accept}
\]
This is exact but mixes poorly in the real-data examples.

\paragraph{Optional tangent proposal.}
For efficiency one may instead propose from
\[
  \PTN\!\left(N_h+\beta,\Omega,B_{ih}+R'_{ih}(\alpha)\right).
\]
This is a state-dependent proposal.  It is only exact if the full
Metropolis--Hastings ratio includes the reverse proposal density, including
the PTN normalizing constants.  Without those proposal-density terms it is an
approximation, not an exact MCMC step.

\section{Sampling the Mean Parameters}

After updating \((\delta_i,\alpha_i)\), sample the inverse means for the
active pattern clusters:
\begin{align}
  q_{ic}=1/\lambda_{ic}
  \mid \alpha_i,\delta_i=h,\bm x_i,a_0,\nu
  \sim
  \Ga\!\left(
    a_0+n_{hc}\alpha_i,\,
    \frac{a_0}{\nu}+\alpha_i S_{ihc}
  \right).
  \label{eq:q-update}
\end{align}
This update is useful for the fully Bayesian hyperparameter updates below.
If hyperparameters are fixed at empirical-Bayes estimates, this step is not
needed for posterior pattern probabilities.

\section{Hyperparameter Updates}

The pattern probabilities have the conjugate update
\[
  \bm\pi\mid \bm\delta
  \sim
  \Dir\!\left(p_0+n_0,\ldots,p_{H-1}+n_{H-1}\right),
  \qquad
  n_h=\sum_{i=1}^G \mathbf{1}\{\delta_i=h\}.
\]
If \(\rho=1/\nu\) and \(\nu\sim \IGa(a_\nu,b_\nu)\), then
\[
  \rho\mid \bm q,a_0
  \sim
  \Ga\!\left(
    a_\nu+a_0M_q,\,
    b_\nu+a_0\sum_{i,c} q_{ic}
  \right),
  \qquad
  \nu=1/\rho,
\]
where \(M_q=\sum_i C_{\delta_i}\) is the number of active gene-pattern
clusters.

If \(\eta=1/\mu\) and \(\mu\sim \IGa(a_\mu,b_\mu)\), then
\[
  \eta\mid \bm\alpha,\beta
  \sim
  \Ga\!\left(
    a_\mu+\beta G,\,
    b_\mu+\beta\sum_{i=1}^G\alpha_i
  \right),
  \qquad
  \mu=1/\eta.
\]

The conditionals for \(a_0\) and \(\beta\) are again gamma-shape targets and
can be updated with the same P-IG--PTN remainder-MH construction.  For example,
up to constants,
\begin{align}
  \log p(a_0\mid \bm q,\nu)
  &=
  (a_{a_0}-1)\log a_0-b_{a_0}a_0
  +M_q a_0\log(a_0/\nu)
  -M_q\log\Gamma(a_0) \nonumber\\
  &\quad
  +(a_0-1)\sum_{i,c}\log q_{ic}
  -\frac{a_0}{\nu}\sum_{i,c}q_{ic}.
  \label{eq:a0-target}
\end{align}
Similarly, up to constants,
\begin{align}
  \log p(\beta\mid \bm\alpha,\mu)
  &=
  (a_\beta-1)\log\beta-b_\beta\beta
  +G\beta\log(\beta/\mu)
  -G\log\Gamma(\beta) \nonumber\\
  &\quad
  +(\beta-1)\sum_{i=1}^G\log\alpha_i
  -\frac{\beta}{\mu}\sum_{i=1}^G\alpha_i.
  \label{eq:beta-target}
\end{align}
Equations \eqref{eq:a0-target} and \eqref{eq:beta-target} contain
\(\Gamma(\cdot)^{-M}\) factors and are therefore natural P-IG updates.

\section{Full MCMC Algorithm}

One MCMC sweep for the \texttt{equalcv=TRUE, nclust=1} GaGa model is:
\begin{enumerate}[label=\arabic*.]
  \item For each gene \(i\), update \(\delta_i\) using
        \eqref{eq:pattern-update}.
  \item For each gene \(i\), update \(\alpha_i\):
        draw \(\tau_c\), draw \(\omega_s\sim\PIG(\alpha_i)\), sample
        \(y_i=\log\alpha_i\) from \eqref{eq:log-alpha-slice} by slice
        sampling, and set \(\alpha_i=\exp(y_i)\).
  \item Sample \(q_{ic}=1/\lambda_{ic}\) for active clusters using
        \eqref{eq:q-update}.
  \item Update \(\bm\pi\) by its Dirichlet conditional.
  \item Update \(\nu\) and \(\mu\) through \(\rho=1/\nu\) and
        \(\eta=1/\mu\).
  \item Update \(a_0\) and \(\beta\) with P-IG remainder-MH steps based on
        \eqref{eq:a0-target} and \eqref{eq:beta-target}, or hold them fixed
        at empirical-Bayes estimates for a direct comparison with
        \texttt{gaga::fitGG}.
\end{enumerate}

Posterior pattern probabilities are estimated by
\[
  \widehat{\Pr}(\delta_i=h\mid\bm x)
  =
  \frac{1}{T}\sum_{t=1}^T
  \mathbf{1}\{\delta_i^{(t)}=h\},
\]
or, with Rao--Blackwellization,
\[
  \widehat{\Pr}(\delta_i=h\mid\bm x)
  =
  \frac{1}{T}\sum_{t=1}^T
  \Pr(\delta_i=h\mid \alpha_i^{(t)},\bm x_i,\Theta^{(t)}).
\]
Differential-expression calls use the same posterior FDR rule as GaGa, with
the all-equal pattern \(h=0\) as the null.

\section{Extensions}

\paragraph{Varying CV, \texttt{equalcv=FALSE}.}
The original GaGa model also permits cluster-specific shapes
\(\alpha_{ic}\).  Conditional on a fixed pattern \(h\), each active cluster
has the same P-IG update as above with \(N_h\) replaced by \(n_{hc}\) and with
only the sufficient statistics \((n_{hc},S_{ihc},L_{ihc})\) for that cluster.
The pattern update then involves a trans-dimensional move unless the
\(\alpha_{ic}\)'s are integrated out.  A practical exact implementation can
use reversible-jump moves, or compute the pattern probabilities by an inner
P-IG sampler for each candidate pattern.

\paragraph{MiGaGa, \texttt{nclust>1}.}
MiGaGa adds a gene-cluster indicator \(d_i\) with component-specific
\((a_{0d},\nu_d)\).  Conditional on \(d_i=d\), all formulas above hold with
\((a_0,\nu)\) replaced by \((a_{0d},\nu_d)\).  The update of \(d_i\) is a
categorical step using the component prior probabilities and the corresponding
collapsed likelihood.

\end{document}
